% Section commune : autoencodeurs sur FPGA (et plateformes voisines)
\section{Autoencodeurs sur FPGA}
\label{sec:autoencodeurs}

La bibliographie de départ supposait que les autoencodeurs (AE) sur FPGA étaient surtout
étudiés sur des séries temporelles et en physique des particules, et peu sur des images.
Les 11 fiches de l'axe permettent de préciser ce constat.

\subsection{Physique des particules et signaux : la très faible latence}

\citet{arxiv-2108.03986-autoencoders-fpgas-real-time-unsupervised} placent des AE et VAE
au déclenchement de premier niveau du LHC, pour sélectionner des collisions anormales sans
hypothèse sur la nouvelle physique. Trois choix sont transposables à d'autres domaines :
\begin{itemize}
  \item le score d'anomalie (erreur quadratique) est calculé sur la puce, par une couche
    ajoutée à hls4ml \src{§IV} ;
  \item le VAE est réduit à son encodeur : le score s'obtient à partir de la divergence de
    Kullback-Leibler ($D_{KL}$) calculée sur $(\mu, \sigma)$, ce qui supprime
    l'échantillonnage et le décodeur \src{§IV} ;
  \item le choix entre QAT et PTQ se fait selon la stabilité du score d'anomalie, et non
    selon la perte d'entraînement \src{§VI}.
\end{itemize}
Le VAE dense (encodeur seul) tient en 80\,ns sur VU9P ; le seul AE convolutif, sur une
entrée de 19$\times$3 valeurs, prend 47\,\% des LUT et est rejeté \src{Table~III}. Ces
chiffres sont des estimations de Vivado HLS ; une seule validation sur carte est faite, et
elle est qualitative \src{§VII}.

Deux autres travaux restent sur des signaux :
\begin{itemize}
  \item \citet{arxiv-2602.15751-enabling-low-latency-machine-learning} compressent des
    impulsions de calorimètre (32 échantillons en 2 valeurs) sur un PolarFire durci : 25\,ns
    de latence en version \enquote{nano}, sans bloc mathématique \src{§5.2}$^{m}$. La
    puissance est estimée par l'outil ;
  \item \citet{2024-que-low-latency-variational-autoencoder-fpgas} sont les seuls à
    échantillonner en matériel (générateur Mersenne Twister suivi d'une transformée de
    Box-Muller), avec un latent de dimension~1, sur des séries d'ondes gravitationnelles
    \src{§III}. Une inférence dure 17 cycles (85\,ns) ; les \enquote{5,6\,ns} mis en avant
    sont un débit inverse amorti sur 10\,000 inférences \src{§IV-E}. La puissance est
    estimée par Vivado \src{§IV-E}.
\end{itemize}
Enfin, \citet{2024-campos-end-end-codesign-hessian-aware} quantifient un encodeur de
compression de données de calorimètre (ECON-AE), dont seul l'encodeur est synthétisé
\src{§7, Table~3}.

\subsection{Images : compression apprise et encodeurs}

\paragraph{Compression apprise sur grosse FPGA.} F-LIC
\citep{2022-sun-f-lic-fpga-learned-image} implante un codec d'images appris (VAE avec
hyper-prior) en virgule fixe 8~bits sur VCU118. Il encode du 720p à 40,69\,FPS et du 1080p
à 19,15\,FPS, avec un PSNR de 29,94\,dB à 0,2542\,bpp \src{Fig.~6}$^{m}$. Deux techniques
sont réutilisables pour un décodeur :
\begin{itemize}
  \item deux poids partagent un même DSP (\enquote{casDSP}) ;
  \item la déconvolution saute les zéros insérés par le stride.
\end{itemize}
Deux limites : le codage entropique tourne sur le PC hôte \src{Fig.~1}, et l'encodeur 720p
occupe 4\,146 DSP. L'énergie par image est comparable à celle d'une Jetson Xavier NX (1,41
contre 1,43\,J en encodage 720p), pas meilleure, et la méthode de mesure n'est pas décrite
\src{Fig.~6}.

\paragraph{Compression à bord, sur VPU.} La mission Φ-Sat-2
\citep{2023-guerrisi-artificial-intelligence-board-image-compression} compresse les images
par un AE convolutif d'un facteur 8 fixe, sans codage entropique \src{§II-A}. L'encodeur
tourne à bord sur une VPU Myriad~2, pas sur un FPGA : 10,92\,s par acquisition pour
1,4\,W, contre 0,95\,s sur un GPU GTX~1650 \src{Table~V}$^{m}$. Le passage à la VPU ne
coûte presque rien en qualité : SSIM de 0,9647 contre 0,9651 \src{§III-B--C}.

\paragraph{Encodeurs d'images sur petite FPGA.}
\citet{2021-manjunath-energy-efficient-edgeai-autoencoder-accelerator} ternarisent
l'encodeur convolutif d'un agent d'apprentissage par renforcement (entrée 120$\times$120
RGB, simulateur) et l'implantent en Verilog sur Artix-7. En 32~bits et 64~PE, une image
coûte 0,8\,ms et 0,25\,mJ \src{Table~5}. La puissance est estimée par Vivado
\src{§VIII-A}, et le décodeur n'est pas implanté.
\citet{arxiv-2301.07050-energy-efficient-reconfigurable-autoencoder-implementation}\pp{}
débruitent MNIST sur XCZU7EV (2,91\,ms, 5,93\,W \enquote{on-chip}). C'est la seule ligne
de \texttt{benchmarks.csv} en confiance \emph{low} : l'\enquote{accuracy} d'un débruiteur
n'est pas définie, et le tableau comparatif est physiquement incohérent \src{Table~II}.

\subsection{Détection d'anomalies sur petite FPGA : des signaux, pas des images}

\begin{itemize}
  \item \citet{2022-malviya-edge-compatible-convolutional-autoencoder-implemented}
    détectent des anomalies de vibrations d'hélicoptère, encodées en scalogrammes
    64$\times$64, avec un AE convolutif de 8\,755 paramètres sur PYNQ-Z2. Encodeur et
    décodeur sont dans deux bitstreams séparés (40\,403 et 38\,169 LUT), pour 2,1 + 3,8\,ms
    \src{Table~4, §IV-B}$^{m}$. La précision est remesurée sur le FPGA et baisse (85,02\,\%
    contre 88,8\,\% sur CPU) en virgule fixe 23~bits \src{§IV-B}.
  \item \citet{2019-tsukada-neural-network-device-learning-anomaly} (ONLAD) apprennent en
    ligne un AE à une couche cachée par OS-ELM (voir
    section~\ref{sec:cl-fpga}). Leurs données sont tabulaires ou des images aplaties.
  \item \citet{2019-coutinho-deep-neural-network-hardware-implementation} : malgré le
    titre, l'\enquote{autoencodeur empilé} n'est qu'un classifieur MLP sur MNIST, sans
    reconstruction \src{§III-A, §V-D}. Leur comparaison met une puissance dynamique de
    quelques mW face à des puissances totales \src{Tables~16--17}.
\end{itemize}

\subsection{Bilan}

\begin{table}[!htbp]
\centering\footnotesize
\caption{Autoencodeurs du corpus : donnée traitée, partie implantée et type de mesure.}
\label{tab:ae}
\begin{tabularx}{\textwidth}{@{}P{2.6cm}P{3.1cm}LP{2cm}P{2.4cm}@{}}
\toprule
Travail & Donnée & Partie en matériel & Puce & Puissance \\
\midrule
AE du LHC\,\textsuperscript{a} & tabulaire 19$\times$3 & AE complet ; VAE encodeur seul & VU9P & non mesurée \\
PicoCal\,\textsuperscript{b} & impulsions 1D & encodeur & PolarFire & estimée \\
VAE de Que\,\textsuperscript{c} & séries temporelles & VAE complet (latent 1) & U250, Z7045 & estimée \\
ECON-AE\,\textsuperscript{d} & calorimètre & encodeur & \nr{} (voir fiche) & non mesurée \\
F-LIC\,\textsuperscript{e} & images 720p et 1080p & VAE hyper-prior, sans codage entropique & VCU118 & méthode non décrite \\
Φ-Sat-2\,\textsuperscript{f} & images satellite & encodeur & VPU (pas de FPGA) & méthode non décrite \\
Manjunath\,\textsuperscript{g} & images 120$\times$120 (simulateur) & encodeur & Artix-7 & estimée \\
CAE débruiteur\,\textsuperscript{h} & MNIST bruité & AE complet & XCZU7EV & méthode non décrite \\
Malviya\,\textsuperscript{i} & vibrations en scalogrammes & AE, 2 bitstreams & XC7Z020 & méthode non décrite \\
ONLAD\,\textsuperscript{j} & tabulaire, MNIST aplati & AE OS-ELM + apprentissage & XC7Z020 & carte (3,1\,W) \\
SSAE\,\textsuperscript{k} & MNIST & MLP (encodeurs empilés) & XC6VLX240T & dynamique \\
\bottomrule
\end{tabularx}

\smallskip\raggedright\footnotesize
\textsuperscript{a}\,\citet{arxiv-2108.03986-autoencoders-fpgas-real-time-unsupervised} ;
\textsuperscript{b}\,\citet{arxiv-2602.15751-enabling-low-latency-machine-learning} ;
\textsuperscript{c}\,\citet{2024-que-low-latency-variational-autoencoder-fpgas} ;
\textsuperscript{d}\,\citet{2024-campos-end-end-codesign-hessian-aware} ;
\textsuperscript{e}\,\citet{2022-sun-f-lic-fpga-learned-image} ;
\textsuperscript{f}\,\citet{2023-guerrisi-artificial-intelligence-board-image-compression} ;
\textsuperscript{g}\,\citet{2021-manjunath-energy-efficient-edgeai-autoencoder-accelerator} ;
\textsuperscript{h}\,\citet{arxiv-2301.07050-energy-efficient-reconfigurable-autoencoder-implementation}\pp ;
\textsuperscript{i}\,\citet{2022-malviya-edge-compatible-convolutional-autoencoder-implemented} ;
\textsuperscript{j}\,\citet{2019-tsukada-neural-network-device-learning-anomaly} ;
\textsuperscript{k}\,\citet{2019-coutinho-deep-neural-network-hardware-implementation}.
\end{table}

Trois constats se dégagent du tableau~\ref{tab:ae} :
\begin{itemize}
  \item \textbf{l'image est traitée}, mais par la compression sur grosse FPGA (F-LIC), sur
    VPU (Φ-Sat-2), ou à petite échelle sur petite FPGA (Manjunath,
    arxiv-2301.07050\pp) ;
  \item \textbf{aucun article ne détecte d'anomalies visuelles par AE sur FPGA} : les AE
    d'anomalie portent sur des signaux ou des données tabulaires ;
  \item \textbf{l'AE complet sur une petite puce reste coûteux} : au 26/09, on implantait
    l'encodeur seul, ou l'on coupait le réseau en deux bitstreams. Depuis, deux articles
    implantent encodeur et décodeur sur une UltraScale+ (CAE d'artefacts hyperspectraux sur
    la DPU d'une ZCU104, \citealp{arxiv-2407.17647-energy-efficient-artefact-detection-accelerator} ;
    compression apprise sur ZCU102,
    \citealp{arxiv-2503.04832-lightweight-embedded-fpga-deployment-learned}). Le coût des
    déconvolutions sur une puce de la taille d'une KV260 reste non mesuré.
\end{itemize}
Le survey d'anomalie visuelle industrielle de
\citet{2023-cui-survey-unsupervised-anomaly-detection-algorithms}, pas encore fiché, dira
quels modèles d'anomalie visuelle seraient les candidats à porter.
